\chapter{Conclusion, Limitations and Future Work}
\label{ch:conclusion}

\section{Summary of Contributions}

This report investigated the architectural security of Lower-Layer Triggered Mobility (LTM) in
5G-Advanced and emerging 6G systems. Rather than focusing on unauthenticated MAC Control Element payloads,
this research evaluated the vulnerabilities created by LTM's core performance optimizations: pre-cached
candidate state, unauthenticated asynchronous timing advance acquisition, and RACH-less switching. Two
fundamental findings were identified through systematic 3GPP specification analysis and validated via
open-source simulation:

\begin{enumerate}\itemsep3pt
  \item \textbf{Finding 1 --- Corrupted Early Timing Advance (\cref{ch:m1}):} Early \TA{} acquisition
        records physical-layer preamble measurements that are bound only to an occasion and preamble
        index, without binding to a UE identity. Furthermore, 3GPP standards impose no retention limits
        on this network-cached value. An adversary using a low-power RF relay can forward preambles with
        excess delay. Measured on srsRAN and OAI, the decodability boundary spans 6--8 units (median 7):
        deterministic uplink failure under flat-fading/AWGN channel conditions by 8 units, with realistic
        multipath saturating slightly more gradually across the same range. As derived in
        \cref{eq:cp9}, this failure boundary is invariant across all 5G and 6G numerologies. The relay
        requires only a $+9.9$\,dB power advantage under flat fading
        ($+5.3$\,dB under multipath), operating from any location closer to the candidate cell than the
        victim.
  \item \textbf{Finding 2 --- Scheduling Request Deadlock (\cref{ch:m2}):} In TS~38.321 clause~5.4.4,
        exhaustion of the Scheduling Request counter (\texttt{SR\_COUNTER}) triggers unconditional
        deallocation of PUCCH, SRS, and all uplink/downlink grants at nesting depth \texttt{5>}. The recovery
        action (initiating Random Access) sits at depth \texttt{6>}, gated by a condition that evaluates
        to false during a RACH-less LTM switch. Formal evaluation across all 5400 legal 3GPP configurations
        shows that this deadlock is reachable in up to 77.4\% of configurations, silencing the UE until
        supervision timer T304 expires.
  \item \textbf{Deterministic Vulnerability Trigger:} Finding~1's corrupted \TA{} is what feeds Finding~2's
        deadlock: the same misalignment that fails the uplink transport block also decorrelates the PUCCH
        Format~0 Scheduling Request (\cref{ch:m2}), driving counter exhaustion with near certainty under
        flat fading and with high probability under multipath. The one escape clause~5.4.4 leaves
        open---a network-issued dynamic grant---does not save the UE, because the same corrupted \TA{}
        that caused the failure makes that grant undecodable too (\cref{sec:m2dyngrant}).
  \item \textbf{Safe Configuration Rule:} We derived the closed-form configuration rule
        $\texttt{sr-TransMax} \times T_{\text{interval}} \ge \texttt{t304}$ that ensures T304 expires
        prior to counter exhaustion, avoiding deadlock at the expense of increased scheduling latency.
\end{enumerate}

\section{Limitations of the Study}
\label{sec:limits}

The theoretical and specification-level findings in this report rest on normative 3GPP specifications
checked clause-by-clause. Several experimental limitations apply to the current phase of work:
\begin{enumerate}\itemsep2pt
  \item \textbf{Simulation-Based Evaluation:} Results are established using open-source protocol stacks
        (srsRAN, OpenAirInterface) and channel emulation models (3GPP TR~38.901). Live over-the-air
        testing on commercial hardware or custom RF relay prototypes was not performed in this phase.
  \item \textbf{Decoupled Stage Modeling:} Physical-layer preamble capture, transport block decodability,
        and Scheduling Request deadlock were evaluated using independent, targeted simulations across
        open-source stacks (srsRAN and OAI). A fully integrated, live multi-node cellular testbed running
        the relay, candidate gNB-DU, and UE MAC stack end-to-end remains an objective for Phase~II.
  \item \textbf{Channel and Receiver Assumptions:} Uplink baseline evaluations utilized single-antenna
        data reception and slot-static channel conditions, although receive diversity was fully swept for
        the control channel.
  \item \textbf{Receiver-Implementation Sensitivity:} The 6--8 unit decodability boundary is measured
        against srsRAN's and OpenAirInterface's specific receiver implementations (channel estimation,
        equalization, and tracking-loop design). Prior work on serving-cell TA manipulation against
        commercial UEs reports a substantially larger RLF-triggering error of 20--60 units
        (\cref{sec:lit-ta}), though that work measures a different failure mode (sustained downlink
        out-of-sync counting) against a different target (the active serving-cell link, under closed-loop
        tracking) than the one-shot coded-BLER cliff measured here against a freshly-acquired candidate
        \TA{}. Whether a production gNB's receiver chain would shift this report's specific 6--8 unit
        boundary is not established in this phase and is left open for further study.
  \item \textbf{Deployment Verification Scope:} Commercial verification was performed against published
        configuration dumps from one vendor (Amarisoft Callbox); wider multi-operator configuration surveys
        remain open.
\end{enumerate}

\section{Future Work}
\label{sec:future}

Phase II will focus on researching and evaluating, through further specification analysis and simulation,
defense mechanisms against the vulnerabilities identified in this phase. The urgency of this work only
grows with 6G numerologies: by \cref{eq:cp9}, the 7-unit failure boundary scales with $T_{\mathrm{unit}}$
and is invariant across subcarrier spacings as an algebraic property of 3GPP's own definitions of
$T_{\mathrm{unit}}$ and $\Delta f$. In upper mid-band (7--24\,GHz) and sub-THz 6G bands, $T_{\mathrm{unit}}$
shrinks from 260\,ns down to 65\,ns ($\mu=3$), compressing the absolute physical time budget available to
absorb a corrupted \TA{} from $1.82\,\mu\text{s}$ to $0.46\,\mu\text{s}$---so any remediation studied now
should be re-examined against that narrower margin for 6G.

\begin{enumerate}\itemsep3pt
  \item \textbf{Protocol \& Specification Hardening:} Refine the TS~38.321 clause~5.4.4 remediation already
        derived in this phase from the clause~5.2 CLTM precedent (\cref{sec:m2gap})---re-aligning resource
        deallocation with recovery actions so that a UE undergoing RACH-less switching retains a fallback to
        Random Access upon counter exhaustion---and examine analogous hardening of the early \TA{} acquisition
        and retention path (\cref{sec:m1retention}) at the specification level.
  \item \textbf{Defensive Telemetry \& Verification:} Study, through specification analysis and simulation,
        network-side defenses that bound and cross-check early timing advance measurements before they are
        applied (including retention limits on cached candidate state and cross-layer consistency checks),
        and physical-layer detection approaches for the multi-path delay and power-delay anomalies
        characteristic of RF relay forwarding.
\end{enumerate}
